% !TeX root = ../main.tex

\chapter{图网络基本理论}
\section{图概述}
图由顶点和边构成, 用以描述对象之间的关系。形式化地, 图可以表示为 \( {G} = \left( {{V}, {E}}\right) \) , 其中: \( {V} \) 是顶点集合, 表示对象。 \( {E} \) 是边集合, 表示顶点之间的关系。

邻居：在图中, 若两个顶点之间存在一条边, 则称它们为邻居。

例如, 在无向图 \( G = \left( {V, E}\right) \) 中, 若存在边 \( e = \left\{  {{v}_i, {v}_j}\right\} \) , 则顶点 \( {v}_i \) , 和 \( {v}_j \) 是邻居。

邻居集合：对于顶点 \( v \) , 其邻居集合 \( N\left( v\right) \) 定义为： \( N\left( v\right)  = \{ u \in  V \mid  \{ v, u\}  \in \)
E \}。

度: 顶点的度是与该顶点直接相连的边的数量。在无向图中, 顶点 \( {v} \) 的度 \( D\left( {v}\right) \)
为: \( D\left( {v}\right)  = \left| {{N}\left( {v}\right) }\right| \)。

在有向图中, 顶点的度分为入度与出度。入度是以该顶点为终点的边的数量。出度是以该顶点为起点的边的数量。

$k$ 阶子图是包含 \(k\) 个顶点及其之间所有可能边的子图。例如，当 \(k = 2\) 时，$k$ 阶子图是由两个顶点和它们之间的边构成的最小子图。当 \(k = 3\) 时，$k$ 阶子图可能是一个三角形，包含三个顶点和三条边。$k$ 阶子图常用于分析网络中的局部结构，如社交网络中的三角形闭包现象。

路径与连通性：路径是从一个顶点到另一个顶点的边序列。路径可以是简单路径（不重复顶点）或回路（起点与终点相同的路径）。若在图中任意两个顶点间都存在路径，则图是连通的。

邻接矩阵：邻接矩阵是描述图的一种矩阵表示形式。对于一个包含 \(n\) 个顶点的图，其邻接矩阵 $A$ 是一个 \(n\times n\) 的矩阵，满足以下条件：若顶点 \(i\) 和顶点 \(j\) 之间存在边，则 \(A[i][j]=1\)（或权值 \(w\)）；若顶点 \(i\) 和顶点 \(j\) 之间不存在边，则 \(A[i][j]=0\)。在无向图中，邻接矩阵是对称的；在有向图中，矩阵的对称性不一定成立。


\section{图数据应用场景}
交通网络：图数据是一种基于节点和边的结构化数据形式，广泛应用于需要描述复杂关系的场景。交通网络是图数据的典型应用之一。在交通网络中，城市中的道路、铁路、公交线路可以表示为图中的边，交叉路口、车站等可作为节点。通过图数据分析，可以实现最短路径规划、交通流量优化、事故风险评估等功能。例如，导航软件利用图算法快速找到最优路径，实时避开拥堵区域，从而提升出行效率。此外，公共交通系统中可以通过图数据优化线路布局，提升整体运输能力。

社交网络：在社交网络中，用户和用户之间的关系可以建模为图，节点表示用户，边表示好友关系或互动行为。通过图算法可以挖掘关键节点（如意见领袖），分析信息传播路径和社区结构。典型应用包括推荐好友、反欺诈检测和舆情分析等。

知识图谱构建：知识图谱通过图结构将概念、实体及其关系进行建模，广泛应用于搜索引擎、智能问答和语义分析中。它可以将自然语言中的信息转化为可查询的结构化数据，帮助实现精准检索、上下文理解和多跳推理。

供应链管理：在供应链中，企业、仓库、运输路径可以建模为图，通过分析节点和边的权重，优化物流路径、降低运输成本。基于图数据，可以提前发现瓶颈环节，保障供应链的稳定性，提高资源利用效率。


\section{图拉普拉斯矩阵}

图信号处理是一门结合图论与信号处理的新兴学科, 用于研究定义在图结构上的信号。图的拉普拉斯矩阵是分析图信号的重要工具, 具有丰富的数学性质和实际意义。本节将简要描述图的拉普拉斯矩阵的定义、形式及其与拉普拉斯算子的关系, 并论述其在反映图信号局部平滑度及构建二次型中的重要作用。

图的拉普拉斯矩阵是描述图结构的重要工具, 广泛应用于图信号处理、网络分析和机器学习中。对于一个无向图 \(G = \left( {V, E}\right)\) , 其拉普拉斯矩阵定义为:
\begin{equation}
L = D - {A}, 
\end{equation}
其中, \(D\) 是度矩阵 (Degree Matrix) : 一个对角矩阵, 其对角线元素 \(D\left\lbrack  {i, i}\right\rbrack\) 表示顶点 \(i\) 的度数。A 是邻接矩阵(Adjacency Matrix):矩阵的元素 \({A}\left\lbrack  {i, j}\right\rbrack\) 表示顶点 \(i\) 和 \(j\) 是否相连。

拉普拉斯矩阵根据其形式可以分为标准型和正则型。

标准型拉普拉斯矩阵 \(L\) 定义为 \(L = D - {A}\)。其元素形式为: 若 \(i = j\) , \(L\left\lbrack  {i, j}\right\rbrack   = D(i)\) (顶点的度数)。

若 \(i \neq  j\) 且 \(i\) 与 \(j\) 相连, \(L\left\lbrack  {i, j}\right\rbrack   =  - 1\)。

若 \(i \neq  j\) 且 \(i\) 与 \(j\) 不相连, \(L\left\lbrack  {i, j}\right\rbrack   = 0\)。

正则型拉普拉斯矩阵是对标准型进行归一化处理的结果, 其定义为:
\begin{equation}
L_\mathrm{norm} = i - {D}^{-1/2}{A}{D}^{-1/2}, 
\end{equation}
这里, I 是单位矩阵, \({D}^{-1/2}\) 是度矩阵的平方根倒数。正则化矩阵通过对顶点的度数进行标准化, 可以更好地平衡图中节点的影响, 特别适用于具有不同度分布的图 \parencite{1997_Chung_Spectral_Graph_Theory}。

拉普拉斯矩阵来源于连续空间中的拉普拉斯算子, 是其在离散图结构上的推广形式。在连续空间中, 拉普拉斯算子用于计算函数的二阶导数, 用于描述局部变化率。其定义为:
\begin{equation}
\Delta {f} = \sum\limits_{{i = 1}}^{{n}}\frac{{\partial }^{2}f}{\partial {x}_i^{2}}.
\end{equation}

在离散图结构中, 拉普拉斯算子退化为拉普拉斯矩阵。具体地, 若定义图信号为一个顶点上的函数值向量 \({f} = {\left\lbrack  {f}_1, {f}_2, \ldots , {f}_{n}\right\rbrack  }^{{T}}\) , 拉普拉斯矩阵作用于图信号时表示为:
\begin{equation}
{{Lf}} = {{Df}} - {{Af}}, 
\end{equation}
这里, \({{Lf}}\) 的结果描述每个顶点的值与其邻居顶点值之间的差异。拉普拉斯矩阵本质上是一个离散的差分算子, 用于反映图上信号的变化程度。

拉普拉斯矩阵是一个反映图信号局部平滑度的算子。对于一个图信号 \({f}\) , 拉普拉斯矩阵可以用来度量信号在图上的平滑程度, 直观地描述信号在邻域中的变化情况。拉普拉斯矩阵可以定义一个非常重要的二次型, 称 \(\sum\limits_{{{e}_{ij} \in  E}}{\left( {x}_i - {x}_j\right) }^{2}\) 为图信号的总变差, 总变差是一个标量, 它将各条边上的信号的差值平方进行加和, 刻画了图信号整体的平滑度。

\begin{equation}
{x}^{T}{Lx} = \sum\limits_{{v}_i}\sum\limits_{{{v}_j \in  N\left( {v}_i\right) }}{x}_i\left( {{x}_i - {x}_j}\right)  = \sum\limits_{{{e}_{ij} \in  E}}{\left( {x}_i - {x}_j\right) }^{2}.
\end{equation}
这一二次型刻画了图信号的平滑性: 当相邻顶点之间的信号值差异较小时, 该值较小；当信号值差异较大时, 该值较大。由此可见, 拉普拉斯矩阵是一个用于分析信号平滑度的重要工具。

可见拉普拉斯矩阵在图信号处理中具有重要意义, 它反映了图信号的局部变化, 能用来衡量图信号的平滑性和不均匀性。它在谱分析中起核心作用, 例如通过特征值分解对图进行聚类或降维。另外拉普拉斯矩阵在图卷积神经网络(GCN)等机器学习模型中被广泛用作构建图滤波器的基础。



\section{图傅里叶变换与图滤波器}

数字信号处理领域中，具有基础性地位的当属傅里叶变换方法。从时域空间向频域空间的转换过程，都是由该变换方法实现的。滤波操作与卷积操作等信号处理技术，其完备的理论定义均建立在该变换基础之上。去噪任务、压缩任务以及重构任务等工程实践应用，也能从中获得明确的理论指导。傅里叶变换同样存在于图信号处理领域，此处将给出相应定义。空域视角到频域视角的转换过程，正是图信号傅里叶变换所要实现的目标，频域视角也将为信号处理带来极大便利。

假设图 \({G}\) 的拉普拉斯矩阵为 \({L} \in  \mathbb{R}^{N * N}\) , 由于 \({L}\) 是一个实对称矩阵, 根据实对称矩阵都可以被正交对角化, 可得:

\begin{equation}
L=V \Lambda V^{T}=\left[\begin{array}{cccc}
\vdots & \vdots & \cdots & \vdots \\
\boldsymbol{v}_1 & \boldsymbol{v}_2 & \cdots & \boldsymbol{v}_N \\
\vdots & \vdots & \cdots & \vdots
\end{array}\right]\left[\begin{array}{lllll}
\lambda_1 & & & \\
& \lambda_2 & & \\
& & \ddots & \\
& & & \lambda_N
\end{array}\right]\left[\begin{array}{ccc}
\cdots & \boldsymbol{v}_1 & \cdots \\
\cdots & \boldsymbol{v}_2 & \cdots \\
\cdots & \vdots & \cdots \\
\cdots & \boldsymbol{v}_N & \cdots
\end{array}\right].
\end{equation}
\({V} \in  {\mathbb{R}}^{{N} \cdot {N}}\) 是一个正交矩阵, 即 \({{{VV}}}^{{T}} = {I}\)。 \({V} = \left\lbrack  {v_{1}, v_{2}, \ldots , v_{{N}}}\right\rbrack\) 表示 \({L}\) 的 \({N}\) 个特征向量, 由于 \({V}\) 是一个正交矩阵, 这些特征向量都是彼此之间线性无关的单位向量。 \(\lambda\) 表示第 \({k}\) 个特征向量对应的特征值, 特征值从小到大排列, 即 \({\lambda }_{1} \leq  {\lambda }_{2}\ldots  \leq  {\lambda }_{{N}}\)。

对于任意给定的向量 \(x\) , 拉普拉斯矩阵的二次型 \({x}^{T}{Lx} = \sum\limits_{{v}_{i}}\sum\limits_{{{v}_{j} \in  N\left( {v}_{i}\right) }}{x}_{i}\left( {{x}_{i} - {x}_{j}}\right)  =\sum\limits_{{{{e}}_{{{ij}}} \in  {E}}}{\left( x_{{i}} - x_{{j}}\right) }^{2}  \geq  0\) , 因此拉普拉斯矩阵是一个半正定型矩阵, 其所有的特征值均大于等于 0。进一步, 可知: \({{LI}} = 0\) , 因此拉普拉斯矩阵具有最小的特征值 0 , 即 \({\lambda }_{1} = 0\)。另外可以证明, 对于 Lsym, 其特征值存在一个上限: \({\lambda }_{{N}} \leq  2\)。

对于任意一个在图 \({G}\) 上的信号 \({x}\) , 其图傅立叶变换为:
\begin{equation}
\widetilde{{x}_{k}} = \sum\limits_{{i = 1}}^{N}{V}_{ki}^{T}{x}_{i} = \left\langle  {{v}_{k}, x}\right\rangle.
\end{equation}
称特征向量为傅里叶基, \(\hat{{x}_{k}}\) 是 \({x}\) 在第 \({k}\) 个傅里叶基上的傅里叶系数。从定义式上可以看到, 傅里叶系数实际上表示了图信号在傅里叶基上的投影，衡量了图信号与傅里叶基之间的匹配程度。通过矩阵运算，可以计算出所有傅里叶系数:
\begin{equation}
\tilde{x} = {V}^{T}x, \tilde{x} \in  \mathbb{R}^{N}.
\end{equation}

由于 \({V}\) 是一个正交矩阵, 对上式左乘 \({V}\) , 则: \({V}\tilde{x} = {V}v^{\prime }{x} = {{Ix}} = {x}\) , 即:
\begin{equation}
{x} = {V}\tilde{x} \in  \mathbb{R}^N.
\end{equation}

于是可以将逆图傅立叶变换定义为:
\begin{equation}
{x}_{k} = \sum\limits_{{i = 1}}^{N}{V}_{ki}{\tilde{x}}_{i}, 
\end{equation}

图滤波器是指对图信号频谱中不同频率成分的强度进行放大或抑制的操作。假设图滤波器为 \({H} \in  {\mathbb{R}}^{{N} * {N}}, {H} : \mathbb{R}^N \rightarrow  \mathbb{R}^N\) , 令输出图信号为 \({y}\) , 则:
\begin{equation}
{y} = {{Hx}} = \sum\limits_{{{k} = 1}}^{{N}}\left( {{h}\left( {\lambda }_{k}\right) {\tilde{x}}_{{k}}}\right) v_{{k}}.
\end{equation}

可以清楚地看到增强还是衰减是通过 \({h}\left( \lambda \right)\) 项来控制的。对上式进行变换:
\begin{equation}
\begin{aligned}
y & =H x=\sum_{k=1}^N\left(h\left(\lambda_k\right) \tilde{x}_k\right) \boldsymbol{v}_k=\left[\begin{array}{cccc}
\vdots & \vdots & \cdots & \vdots \\
\boldsymbol{v}_1 & \boldsymbol{v}_2 & \cdots & \boldsymbol{v}_N \\
\vdots & \vdots & \cdots & \vdots
\end{array}\right]\left[\begin{array}{c}
h\left(\lambda_1\right) \tilde{x}_1 \\
h\left(\lambda_2\right) \tilde{x}_2 \\
\vdots \\
h\left(\lambda_N\right) \tilde{x}_N
\end{array}\right] \\
& =V\left[\begin{array}{cccc}
h\left(\lambda_1\right) & & & \\
& h\left(\lambda_2\right) & & \\
& & \ddots & \\
& & & h\left(\lambda_N\right)
\end{array}\right]\left[\begin{array}{c}
\tilde{x}_1 \\
\tilde{x}_2 \\
\vdots \\
\tilde{x}_N
\end{array}\right] \\
& =V\left[\begin{array}{cccc}
h\left(\lambda_1\right) & & & \\
& h\left(\lambda_2\right) & & \\
& & \ddots & \\
& & & h\left(\lambda_N\right)
\end{array}\right] V^{T} \boldsymbol{x}
\end{aligned}.
\end{equation}

因此, 也可以得到 \({H}\) 的形式:
\begin{equation}
H = V\left\lbrack  \begin{matrix} h\left( {\lambda }_{1}\right) & \cdots & 0 \\  \vdots &  \ddots  & \vdots \\  0 & \cdots & h\left( {\lambda }_{N}\right)  \end{matrix}\right\rbrack  {V}^{T} = V{\Lambda }_{h}{V}^{T}.
\end{equation}

相较于拉普拉斯矩阵, \({H}\) 矩阵只在对角与边坐标上时才有可能取非零值。从运算角度看，$Hx$ 代表在每个节点的局部子图上进行变换。通常，满足此性质的矩阵称为图位移算子，例如，拉普拉斯矩阵和邻接矩阵。

称 \({\Lambda }_{{h}}\) 为图滤波器 \({H}\) 的频率响应矩阵, 对应的函数 \({h}\left( \lambda \right)\) 为 \({H}\) 的频率响应函数, 不同的频率响应函数可以实现不同的滤波效果。在信号处理中, 常见的滤波器有 3 类:低通滤波器、高通滤波器、带通滤波器(见图 \ref{fig:filter})。

\begin{figure}[htbp]
  \begin{center}
    \includegraphics[width=0.9\textwidth]{01960c08-a17d-7162-9acb-7f9eeb4038d6_4_299_1466_1099_255_0.jpg}
    \end{center}
  \caption{三类滤波器的频率响应函数 \parencite{2021_刘忠雨_深入浅出图神经网络}}
  \label{fig:filter}
\end{figure}


低频成分被低通滤波器所保留，信号中平滑部分由此得到凸显; 从 \({b}\) 图中可以看出, 高频成分则由高通滤波器实现选择性保留，快速变化区域成为重点关注对象; 从 \({c}\) 图中可以看出, 特定频段成分的保留通过带通滤波器完成。


\section{图卷积神经网络}

在网络设计领域内，存在着若干关键性思路值得探讨。频率响应矩阵的参数化过程首先需要被关注。图卷积运算与图滤波操作本质相同。作为图滤波算子的核心构成要素，频率响应矩阵的重要地位由此凸显。因此，对该矩阵实施参数化处理具有明显必要性。基于此考量，便可以顺理成章定义如下所示的神经网络层结构：

\begin{equation}
{X}^{\prime } = \sigma \left( {V\left\lbrack  \begin{matrix} {\theta }_{1} & \cdots & 0 \\  \vdots &  \ddots  & \vdots \\  0 & \cdots & {\theta }_{N} \end{matrix}\right\rbrack  {V}^{T}X}\right)  = \sigma \left( {V\operatorname{diag}\left( \theta \right) {V}^{T}X}\right)  = \sigma \left( {\theta X}\right).
\end{equation}

其中 \(\sigma \left( \cdot \right)\) 是激活函数, \(\theta  = \left\lbrack  {{\theta }_{1}, {\theta }_{2}, \ldots , {\theta }_{N}}\right\rbrack\) 是需要学习的参数, \(\odot\) 是对应的需要学习的图滤波器, \({X}\) 是输入的图信号矩阵, \({{X}}^{\prime }\) 是输出的图信号矩阵。

这个设计理念虽然直观且易于理解, 存在显著问题的是，引入的学习参数量级过大，与图中节点数呈现正相关。对于大规模图数据而言，上亿节点的图尤其常见，容易出现过拟合的结果。因此，该网络层设计思路导致学习参数规模产生膨胀化趋势，直接影响了模型的泛化性能。

实际图数据中，集中于低频段的往往是信息的有效部分，高频段的信息承载量呈现出明显的弱化态势。设置多个维度的自由度以进行图滤波器设计，并不合理，有些频率的学习需求并不显著。

\textcite{2016_Defferrard_Convolutional_neural_networks_on_graphs_with_fast_localized_spectral_filtering} 提出了一种新的图卷积神经网络方法, 旨在提高图数据处理的效率和性能。该方法基于图的谱理论, 通过快速局部化的频谱滤波器来实现卷积操作。这些滤波器在频谱域上进行设计, 使得卷积操作不仅能有效地处理图中的节点, 还能捕捉节点之间复杂的结构关系。他们提出了一个新型的图卷积框架, 通过频谱滤波器来代替传统的邻接矩阵操作, 避免了昂贵的计算成本。其主要思想方法是将拉普拉斯矩阵的多项式表示转换为一种可以进行学习的形式，以便更好地拟合不同的频率响应函数。具体如下:
\begin{equation}
{X}^{\prime } = \sigma \left( {V\left( {\mathop{\sum }\limits_{{k = 0}}^{k}{\theta }_{k}{\Lambda }^{k}}\right) {V}^{T}X}\right)  = \sigma \left( {{Vdiag}\left( {\psi \theta }\right) {V}^{T}X}\right), 
\end{equation}
多项式系数向量由 \(\theta = \left\lbrack {{\theta }_{1}, {\theta }_{2}, \ldots , {\theta }_{N}}\right\rbrack\) 构成，作为网络层中真正需要被学习到的参数量化表征。区别于先前论述的方法，参数量 \({K}\) 在该框架下具备可调控特性。随着 \({K}\) 数值的增大，频率响应函数所能拟合的最高次数呈现上升态势，就能够刻画输入图信号矩阵与输出图信号矩阵间存在的复杂滤波关联性。相反地，当 \({K}\) 取值较小时，频率响应函数的拟合能力受限，仅能拟合输入与输出之间较为基础的滤波作用模式。通常，常规情况下设定 \({K} < < {N}\) 即可，可以显著降低模型过拟合情况的发生概率。

在上述两种方法中，虽然网络层的参数量减少，但计算复杂度却因特征分解过程的强依赖性而急剧提升。为了解决这个问题, \textcite{2016_Kipf_Semi_supervised_classification_with_graph_convolutionalnetworks} 提出直接通过邻接矩阵和节点特征的卷积操作来学习节点表示, 对上式进行了限制, 设 \({K} = 1\) , 则:
\begin{equation}
{X}^{\prime } = \sigma \left( {{\theta }_{0}X + {\theta }_{1}{LX}}\right).
\end{equation}
令 \({\theta }_{0} = {\theta }_{1} = \theta\) , 则:
\begin{equation}
{{X}}^{\prime } = \sigma \left( {\theta \left( {{I} + {L}}\right) {X}}\right)  = \sigma \left( {\theta \widetilde{{L}}{X}}\right), 
\end{equation}

在网络学习过程中数值稳定性的提升方面，研究者\textcite{2016_Kipf_Semi_supervised_classification_with_graph_convolutionalnetworks}进行了探索。受到正则拉普拉斯矩阵的启发，归一化处理被应用于L矩阵。定义\(\widetilde{A} = A + I\)的情况下，构建了重归一化形式的拉普拉斯矩阵\({\widetilde{L}}_{sym} = {\widetilde{D}}^{-1/2}\widetilde{A}{\widetilde{D}}^{-1/2}\)这一数学表达式。重归一化形式的拉普拉斯矩阵特征值在(-1, 1)区间内。由此可见，在多层网络训练过程中使用重归一化的拉普拉斯矩阵能够有效抑制梯度消失或爆炸这类优化难题。

为了加强网络的拟合能力, 并设计了一个参数化的权重矩阵 \({W}\) 对输入的图信号矩阵进行仿射变换, 于是得到:
\begin{equation}
{X}^{\prime } = \sigma \left( {{\widetilde{L}}_{sym}{XW}}\right).
\end{equation}

相比于用频率响应函数的拟合形式，图卷积层呈现出显著简化性。最终图卷积相应的图滤波器退化为 \({\widetilde{L}}_{sym}\) 这一形式。转换至运算层面，图卷积操作则演变为 \({\widetilde{L}}_{sym}{X}\) 的矩阵乘积形式。当 \({X}\) 从信号矩阵转换为特征矩阵时，具有图位移算子性质的 \({\widetilde{L}}_{sym}\) 显现出特殊计算特性。依据矩阵乘法的行向量解释方式可见，该算子的运算过程本质上执行着邻居节点特征向量的聚合化操作。


\section{GCN的性质}
GCN 能够对图数据进行端对端学习。近年来深度学习技术取得显著进展之际，作为其核心组成部分的端对端学习范式正日益凸显其重要性。实例表明，深度学习所取得的突破性成果在很大程度上依赖于端对端学习机制的有效运作。端对端的学习方法具备从原始数据中自动提取有效特征的能力。需要指出的是，这种高度自动化学习特性的实现过程往往伴随着大量针对特定数据类型的适配性工作。在当前主流深度学习架构中，各种精心设计的网络层或层块结构便是明证。以二维卷积层为例，其对图像数据的特征提取具有显著优势；而长短期记忆网络层则在处理序列数据时表现出独特效能；全局池化层对于宏观特征的捕获同样不可或缺。由此可见，这些专用计算模块必须严格遵循数据内在规律进行优化设计。实施端对端优化的图卷积网络（GCN），相较于本文基线模型采用的图嵌入技术而言，存在着显著差异：当节点特征向量完成拼接操作后即进入固化状态，下游任务产生的监督信号难以有效作用于图数据表示学习过程，这种机制可能导致节点特征表示无法达到最优任务适配度。相反地，完成图数据建模的GCN并未采用两阶段处理模式，节点的表示学习与后续任务学习被整合到整体模型框架内进行协同优化。在节点表示学习质量的提升方面，属性信息的补充发挥着重要作用，特别是针对那些结构稀疏性较强的图数据而言。由GCN模型实现的是结构信息与属性信息的同步学习过程，结构信息中所蕴含的知识具有独特性，这些知识在属性信息中并不存在，对于节点的准确刻画至关重要。被放置于同一个网络层中进行协同学习的结构信息和属性信息，二者能够共同作用于最终节点表示的形成。由此可见，直接在原始图数据上构建学习过程的GCN模型，为图数据的处理提供了一套端对端的解决方案框架化。在相关任务的学习适应性方面，该模型展现出了明显优势性。

GCN 是一个低通滤波器。回到 GCN 的核心计算式 \({\widetilde{L}}_{sym}{XW}\) 上, 体现图滤波的地方就在于左乘了一个重归一化形式的拉普拉斯矩阵 \({\widetilde{L}}_{sym}\) , 要确定是否为低通滤波, 我们就必须去研究 \({\widetilde{L}}_{sym}\) 对应的频率响应函数 \({p}\left( \lambda \right)\) 的性质。
\begin{equation}
{\widetilde{L}}_{\text{ sym }} = {\widetilde{D}}^{-1/2}\widetilde{A}{\widetilde{D}}^{-1/2} = {\widetilde{D}}^{-1/2}\left( {\widetilde{D} - L}\right) {\widetilde{D}}^{-1/2} = I - {\widetilde{D}}^{-1/2}L{\widetilde{D}}^{-1/2} = I - {\widetilde{L}}_{s}.
\end{equation}

由于 \({\widetilde{L}}_{s}\) 可以被正交对角化, 因此设 \({\widetilde{L}}_{s} = V\widetilde{\Lambda }{V}^{T}\) 是 \({\widetilde{L}}_{s}\) 的特征值, 可以证明, \({\widetilde{\lambda }}_{i} \in  \left\lbrack  {0, 2}\right\rbrack\)。
因此 \({\widetilde{L}}_{sym}\) 可以改写成:
\begin{equation}
{\widetilde{L}}_{\text{ sym }} = I - V\widetilde{\Lambda }{V}^{T} = V\left( {1 - \widetilde{\Lambda }}\right) {V}^{T}.
\end{equation}

显然, 其频率响应函数为 \({p}\left( \lambda \right)  = 1 - {\widetilde{\lambda }}_{i}, \in  ( - 1, 1\rbrack\) , 该函数是一个线性收缩的函数, 因此能起到对图信号进行低通滤波的作用。如果将信号矩阵 \({X}\) 不断左乘 \({K}\) 次 \({\widetilde{L}}_{sym}\) , 则对应频率响应函数为 \({\left( 1 - {\widetilde{\lambda }}_{i}\right) }^{K}\) , 图 \ref{fig:response} 为该函数的图像:


\begin{figure}[htbp]
  \begin{center}
    \includegraphics[width=1\textwidth]{01960c08-a17d-7162-9acb-7f9eeb4038d6_9_220_314_1175_599_0.jpg}
  \end{center}
  \caption{响应函数图像 \parencite{2021_马耀_图深度学习}}
  \label{fig:response}
\end{figure}

从图中可以看到, 随着 \({K}\) 的增大, 频率响应函数在低频段上有着更强的缩放效果, 因此是一种更强效应的低通滤波器。这种堆叠式的滤波操作, 在一定程度上解释了多层 GCN 模型对于信号的平滑能力。事实上, 为了更好地突出这种能力、减少模型的参数量, 有些研究直接将多层 \({{GCN}}\) 退化成 \({o}\left( {{\widetilde{L}}_{sym}{{XW}}}\right)\)。

\textcite{2019_Maehara_Revisiting_Graph_Neural_Networks_All_We_Have_is_Low_Pass_Filters} 提出了一个关于图数据的假设一一输入数据的特征信号包括低频信号与高频信号, 而低频信号包含着对任务学习更加有效的信息。

与此同时, GCN 也存在一定的问题。在使用多层 GCN 之后, 节点的区分性变得越来越差, 节点的表示向量趋于一致, 这使得相应的学习任务变得更加困难。我们将 这个现象称为多层 GCN 的过平滑问题。

GCN 相当于对输入信号做了一个低通滤波的操作, 这会使信号变得更加平滑, 这也是 GCN 模型的一个内在优势。但是, 过犹不及, 多次执行这类对信号不断平滑的操作之后, 信号会越来越趋同, 也就丧失了节点特征的多样性。

GCN 过平滑问题的本质是信息传播次数过多, 导致节点特征趋同, 使得模型的区分能力下降。数学上, 这可以归因于归一化邻接矩阵的特征值收敛。

\begin{equation}
{{H}}^{\left( {l} + 1\right) } = \sigma \left( {\widetilde{{A}}{{H}}^{\left( {l}\right) }{{W}}^{\left( {l}\right) }}\right)
\end{equation}
\begin{equation}
{{H}}^{\left( {L}\right) } = {\left( \widetilde{{A}}\right) }^{{L}}{{H}}^{\left( 0\right) }{{W}}^{\left( 0\right) }\ldots {{W}}^{\left( {L} - 1\right) }.
\end{equation}

如果 \({L}\) 很大, 由于 \(\widetilde{{A}}\) 是归一化的, 它的最大特征值小于 $1$, 因此不断迭代后, 所有节点的特征都会趋于一个稳定的状态, 即所有节点的表示向量变得相似, 从而导致分类能力下降。



\section{其他经典交通流量预测模型}
本节将介绍经典交通流量预测模型, 本章所介绍的模型将作为下一章节的对比模型使用。

\subsection{长短记忆神经网络}

在使用 LSTM 模型 \parencite{2016_Greff_LSTM_A_search_space_odyssey} 预测交通流量时, 虽然 LSTM 主要专注于学习时间依赖性, 但其网络结构中的层次结构也能够在一定程度上捕捉到节点间的关系。LSTM 的每一层都有多个权重参数, 这些参数不仅用于学习时间序列中的依赖性, 还能通过全连接层的线性组合, 捕捉到不同时间步输入之间的相关性。通过这些线性组合, LSTM 能够有效地“学习”到不同特征之间的复杂关系, 这种特性在某些情况下可以间接地反映出交通流量节点之间的潜在关联。

RNN 是一种基础的循环神经网络, 通过循环连接将前一时刻的输出传递给当前时刻, 用于捕捉序列中的时间依赖性。然而, RNN 存在“梯度消失” 问题, 导致它在处理长序列时, 无法有效记住远距离的信息。LSTM 是对 RNN 的一种改进, 设计了 “门”机制来控制信息的传递。LSTM 通过引入输入门、遗忘门和输出门, 能够选择性地保存和忘记信息, 从而有效解决了梯度消失问题, 能够捕捉长期依赖关系。相较于普通的 RNN, LSTM 能够更好地处理长时间间隔的依赖性, 尤其在自然语言处理和语音识别等任务中表现更好。图 \ref{fig:lstm} 为 LSTM 的模型结构图。


\begin{figure}[htbp]
  \begin{center}
    \includegraphics[width=0.6\textwidth]{01960c08-a17d-7162-9acb-7f9eeb4038d6_11_574_316_544_393_0.jpg}
    \end{center}
  \caption{LSTM 模型结构图}
  \label{fig:lstm}
\end{figure}

序列的每一个位置, 输入是当前时刻的 \({x}_{t}\) , 以及信息流中前一时刻的隐藏状态 \({h}_{t - 1}\) 和单元状态 \({c}_{t - 1}\) , 输出是当前时刻的隐藏状态 \({h}_{t}\) 与单元状态 \({c}_{t}\)。

LSTM 包含三个主要的门: 遗忘门 \({{f}}_{{t}}\) 、输入门 \({{i}}_{{t}}\) 、输出门 \({{o}}_{{t}}\) , 以及单元候选状态 \({\widetilde{c}}_{{t}}\) ；

遗忘门 \({f}_{t}\) 决定了前一时刻的单元状态 \({c}_{t - 1}\) 有多少比例被遗忘。
\begin{equation}
{f}_{t} = \sigma \left( {{W}_{f}{x}_{t} + {U}_{f}{h}_{t - 1} + {b}_{f}}\right).
\end{equation}

输入门决定了当前输入 \({x}_{t}\) 和前一时刻的隐藏状态 \({h}_{t - 1}\) 对单元状态的影响。
\begin{equation}
{i}_{t} = \sigma \left( {{W}_{i}{x}_{t} + {U}_{i}{h}_{t - 1} + {b}_{i}}\right).
\end{equation}

候选单元状态是当前时刻的潜在记忆更新。
\begin{equation}
{\widetilde{c}}_{t} = \tanh \left( {{W}_{c}{x}_{t} + {U}_{c}{h}_{t - 1} + {b}_{c}}\right).
\end{equation}

当前的单元状态是由遗忘门调节的上一时刻的状态 \({c}_{t - 1}\) 和输入门调节的候选状态 \({\widetilde{c}}_{t}\) 加权合成的。
\begin{equation}
{c}_{t} = {f}_{t}{c}_{t - 1} + {i}_{t}{c}_{t}.
\end{equation}

当前的隐藏状态是通过输出门调节的单元状态 \({c}_{t}\) 的非线性变换。输出门决定了从当前单元状态 \({c}_{t}\) 中输出多少信息到隐藏状态中。
\begin{equation}
{o}_{t} = \sigma \left( {{W}_{o}{x}_{t} + {U}_{o}{h}_{t - 1} + {b}_{o}}\right)
\end{equation}
\begin{equation}
{h}_{t} = {o}_{t} \odot  \tanh \left( {c}_{t}\right).
\end{equation}

因此, 可以将 LSTM 神经网络模块中用于处理 t 位置的式总结为:
\begin{equation}
{c}_{t}, {h}_{t} = \operatorname{LSTM}\left( {{x}_{t}, {c}_{t - 1}, {h}_{t - 1}}\right).
\end{equation}

LSTM 层的输出通常包括每个时间步的隐状态和单元状态。在进行最终预测时, 可以根据任务的不同, 选取不同的输出方式。对于序列标注任务, 可以使用最后一个时间步的隐状态作为整个序列的表示进行预测; 对于序列生成任务, 则可以逐步使用 LSTM 的输出生成下一步的结果。通过此机制, 使得模型能够在交通流预测中有效地利用过去的状态来预测未来的交通状况 \parencite{2022_刘建伟_循环神经网络研究综述}。此外, LSTM 的输出还可以通过全连接层或其他网络层进一步处理, 映射到所需的预测维度, 例如分类或回归任务的输出。这些输出通常经过 softmax 等激活函数, 转化为最终的预测结果。


\subsection{图嵌入方法}
图嵌入的主要目标是将图中每个节点映射到一个低维向量空间, 从而在嵌入空间中有效保留图的结构信息。这些低维向量表示通常被称为节点嵌入, 能够反映图中节点的某些关键信息。对于哪些信息应当被保留, 已有研究探讨了多种信息类型的保留方式, 包括节点的邻域结构信息、节点的结构角色、社区信息以及节点的状态特征等。至于如何有效地保留这些信息, 近年来提出的多种技术方案对如何保留这些信息做出了回应。尽管具体实现细节各异, 但大多数方法的共同理念是利用嵌入空间中的节点表示来重建原图中的关键信息。其主要动因在于, 理想的节点嵌入应具备有效重建目标信息的能力。
当一些节点在随机游走过程中经常共同出现时, 这些节点被认为是相似的。为了捕捉这种相似性, 通常会优化一个映射函数, 使得学到的节点表示能够有效重构从随机游走中提取的节点间共现信息。DeepWalk \parencite{2014_Perozzi_DeepWalk_Online_Learning_of_Social_Representations} 是一种经典的图嵌入算法, 用于保留节点间的共现关系。本节首先概述了 DeepWalk 算法的整体框架, 并详细阐述了其映射函数、特征提取器、重构器以及目标函数。随后, 将介绍本文重点使用的保留节点共现关系的算法——node2vec.
定义映射函数 \(f\left( {v}_{i}\right)\) 最简洁、方便的方式即使用查找表, 即 \(f\left( {v}_{i}\right)  = {u}_{i} = {W}^{T}{e}_{i}\) .

设 \({G} = \{ {V}, {E}\}\) 表示一个连通图。考虑从图 \({G}\) 上的节点 \({v}\left( 0\right)  \in  {V}\) 开始的随机游走。假设在随机游走的第 \({t}\) 步访问的是节点 \({v}\left( {t}\right)\) , 那么随机游走的下一个节点按照如下概率从 \({v}\left( {t}\right)\) 的邻居节点中选出:
\begin{equation}
p\left( {v\left( {t + 1}\right)  \mid  v\left( t\right) }\right)  = \left\{  \begin{matrix} 1/d\left( {v\left( t\right) }\right), & v\left( {t + 1}\right)  \in  N\left( {v\left( t\right) }\right) \\  0, & \text{ 其他 } \end{matrix}\right.
\end{equation}

两个节点间的共现被标示为元组 \(\left( {{{v}}_{{{con}}}, {{v}}_{{{cen}}}}\right)\) , 其中 \({{v}}_{{{con}}}\) 表示上下文节点, \({{v}}_{{{cen}}}\) 表示中心节点。

为了生成随机游走能够捕获整个图的信息, 每个节点都被用作起始节点生成 \(\gamma\) 个随机游走。因此, 在遍历整个图后, 总共会获得 \(\left| V\right|  \cdot  \gamma\) 个随机游走。对于每个随机游走, 遍历其中的节点, 并将所有的 \(\left( {{{v}}^{{i} - {j}}, {{v}}^{{i}}}\right)\) 和 \(\left( {{{v}}^{{i} + {j}}, {{v}}^{{i}}}\right)\) 添加到共现列表 I 中。

通俗来说, 节点嵌入的主要思想即为: 两个节点在共现列表 I 中出现越多, 说明两个节点的关系约密切, 节点的结构信息越相像, 二者嵌入生成的向量在对方的嵌入向量方向上的投影越大。在这里引入节点扮演不同角色时的不同向量表示: 中心节点与上下文节点, 即 \({{f}}_{{{cen}}}\left( {{v}}_{{i}}\right)  = {{u}}_{{i}} = {{W}}_{{{cen}}}^{{T}}{{e}}_{{i}}, {{f}}_{{{con}}}\left( {{v}}_{{i}}\right)  = {{v}}_{{i}} = {{W}}_{{{con}}}^{{T}}{{e}}_{{i}}\) .

两个节点的共现关系可以表述为: 在中心节点 \({A}\) 的上下文中观测到了 \({B}\)。故对全部的共现关系 I 进行概率建模:
\begin{equation}
p\left( {{v}_{\text{ con }} \mid  {v}_{\text{ cen }}}\right)  = \frac{\exp \left( {{f}_{\text{ con }}{\left( {v}_{\text{ con }}\right) }^{T}{f}_{\text{ cen }}\left( {v}_{\text{ cen }}\right) }\right) }{\mathop{\sum }\limits_{{v \in  V}}\exp \left( {{f}_{\text{ con }}{\left( v\right) }^{T}{f}_{\text{ cen }}\left( {v}_{\text{ cen }}\right) }\right) }.
\end{equation}

一个被良好重构向量表示是能够准确地推断出 I 的共现信息的。为了实现这一目标, 应该为从 I 中提取的元组输出高的概率值, 该思想类似于极大似然估计的思想, 选择那些使得已观察到的数据在给定模型下出现的概率最大的参数值。随机游走假设共现 I 中的这些元组是相互独立的。因此, 重构 I' 的概率可以建模如下, 使得下式的取最大:
\begin{equation}
{{I}}^{\prime } = \mathop{\prod }\limits_{\left( \left( {{v}}_{{{con}}}, {{v}}_{{{cen}}}\right)  \in  {I}\right) }{p}\left( {{{v}}_{{{con}}} \mid  {{v}}_{{{cen}}}}\right).
\end{equation}

node2vec 引入了一种更灵活的方式生成随机游走, 它通过有偏的随机游走探索给定节点的邻域, 并用它代替 DeepWalk 中的随机游走来生成工。具体地说, node2vec 提出了具有两个超参数 \({p}\) 和 \({q}\) 的二阶随机游走, 其定义如下。
\begin{equation}
a\left( {v\left( {t + 1}\right)  \mid  v\left( {t - 1}\right) , v\left( t\right) }\right)  = \left\{  \begin{matrix} 1/p\operatorname{dis}\left( {v\left( {t - 1}\right) , v\left( {t + 1}\right) }\right)  = 0 \\  1\operatorname{dis}\left( {v\left( {t - 1}\right) , v\left( {t + 1}\right) }\right)  = 1 \\  1/q\operatorname{dis}\left( {v\left( {t - 1}\right) , v\left( {t + 1}\right) }\right)  = 2 \end{matrix}\right. .
\end{equation}

该算法由 \textcite{2016_grover_node2vec} 提出。Node2Vec 通过参数化的随机游走策略来调整节点之间的关系。具体来说, Node2Vec 通过两个参数 \({p}\) 和 \({q}\) 控制随机游走的转移概率, 决定了在游走中如何平衡局部和全局结构:

参数 \({p}\) 控制从当前节点返回到上一节点的概率。较小的 \({p}\) 值会增加返回到上一节点的概率, 进而使得游走更加局部化, 专注于节点的局部邻域。参数 \({q}\) 控制游走是否偏向于访问远离当前节点的节点。较小的 \({q}\) 值会导致游走更加探索邻近节点, 类似于 BFS （广度优先搜索）; 较大的 $q$ 值会使游走偏向于访问较远的节点, 类似于 DFS （深度优先搜索）。这种灵活的随机游走策略允许算法能够同时捕捉到局部邻域和全局结构特征。

由于城市数据（例如交通条件）的演变受到底层道路网的限制, 因此将道路网信息纳入预测模型至关重要。为此, 相关文献提出了一种空间嵌入来将顶点编码为向量, 以保留图的结构信息。顶点特征的学习过程通过node2vec方法实现。要将预学习所得向量与整体模型进行协同训练，则需将这些向量输入至包含两个全连接层的神经网络结构之中进行维度上的对齐，最终形成空间嵌入。表示为 \({{e}}_{{{v}}_{{i}}}^{{S}} \in  {\mathbb{R}}^{{D}}\), 其中 \({{v}}_{{i}} \in  {v}\) .

\textcite{2020_Zheng_GMAN_A_graph_multi_attention_network_for_traffic_prediction} 进一步提出了一个时间嵌入, 以编码每个时间步长到一个向量。将空间嵌入与时间嵌入进行维度上的对齐后，融合为时空嵌入，并联合节点特征，使用转换注意力机制进行预测。



\subsection{T - GCN}

该模型由 \textcite{2019_Zhao_T_gcn_A_temporal_graph_convolutional_network_for_traffic_prediction} 在 DCRNN 模型的基础上, 提出一种基于 GCN 和 GRU 的时间图卷积网络模型, T-GCN 模型是为了解决交通预测问题而提出的一种新型深度学习框架。该模型在传统的图卷积网络（GCN）基础上引入了时间序列的建模能力, 能够同时捕捉交通流数据中的时序变化和图结构的空间依赖关系。

在门控递归单元与图卷积网络相融合的基础上，构建了T-GCN模型架构。空间依赖性的建模，是通过图卷积网络捕获路网拓扑结构完成的。交通数据时间维度上的动态变化特征，则使用门控递归单元进行学习。

具体而言, 即每个时间片段中单独使用 GCN 学习结合网络拓扑结构后的表示, 然后再将这些表示作为时序数据输入到循环神经网络中, 通过单元间的信息传递获得动态变化, 捕捉时间特征。最后, 通过全连通层得到结果。模型结构图如图 \ref{fig:t-gcn} 所示。
\begin{figure}[htbp]
  \begin{center}
    \includegraphics[width=0.6\textwidth]{01960c08-a17d-7162-9acb-7f9eeb4038d6_15_462_957_711_310_0.jpg}
    \end{center}
  \caption{T-GCN 模型结构图 \parencite{2023_胡晓荣_基于图神经网络的路网短时交通流预测算法研究}}
  \label{fig:t-gcn}
\end{figure}